\documentclass[conference]{IEEEtran}
\IEEEoverridecommandlockouts
\usepackage[T1]{fontenc}
\usepackage{cite}
\usepackage{amsmath,amssymb,amsfonts}
\usepackage{graphicx}
\usepackage{textcomp}
\usepackage{xcolor}
\usepackage{array}
\usepackage{tabularx}
\usepackage[hyphens]{url}
\renewcommand{\tabularxcolumn}[1]{>{\raggedright\arraybackslash}p{#1}}
\begin{document}

\title{Evaluating the Transport Impact of Large\\Post-Quantum DNSSEC Signatures\\
{\large Theoretical Part (Project -- Part~1)}
\thanks{Network and Information Systems Security, ISCTE -- Instituto
Universit\'ario de Lisboa. Submitted on 28~October~2026.}
}

\author{\IEEEauthorblockN{Justin Kasper}
\IEEEauthorblockA{\textit{Student no.\ 145971}\\
\textit{ISCTE -- Instituto Universit\'ario de Lisboa}\\
Lisbon, Portugal}
\and
\IEEEauthorblockN{Mert Duman}
\IEEEauthorblockA{\textit{Student no.\ 145872}\\
\textit{ISCTE -- Instituto Universit\'ario de Lisboa}\\
Lisbon, Portugal}
\and
\IEEEauthorblockN{Efim Shchegolev}
\IEEEauthorblockA{\textit{Student no.\ 144376}\\
\textit{ISCTE -- Instituto Universit\'ario de Lisboa}\\
Lisbon, Portugal}
\and
\IEEEauthorblockN{Sezer Tetik}
\IEEEauthorblockA{\textit{Student no.\ 145873}\\
\textit{ISCTE -- Instituto Universit\'ario de Lisboa}\\
Lisbon, Portugal}
}

\maketitle

\begin{abstract}
Post-quantum signature schemes such as ML-DSA-44, standardized by NIST in
FIPS~204, produce signatures and public keys that are one to two orders of
magnitude larger than those of the classical algorithms used in DNSSEC today.
As a result, signed DNS responses routinely exceed the 1,232-byte EDNS(0)
payload recommendation and the typical Internet MTU, which pushes them onto IP
fragmentation, truncation with TCP fallback, or alternative transports. This
theoretical part reviews the DNS and DNSSEC foundations, the classical and
post-quantum signature algorithms relevant to DNSSEC, and the message-size and
transport mechanisms that govern how large responses are delivered, including
TCP fallback, Application-layer Request-based Resource Record Fragmentation
(ARRF), and DNS over QUIC (DoQ). It then surveys related work and identifies
the gap addressed by the practical part of the project: a controlled,
reproducible comparison of these transports for ML-DSA-44 against an ECDSA
baseline under varying latency and packet loss.
\end{abstract}

\begin{IEEEkeywords}
DNSSEC, post-quantum cryptography, ML-DSA, TCP fallback, IP fragmentation,
ARRF, DNS over QUIC
\end{IEEEkeywords}

\section{Introduction}
\label{sec:intro}

The Domain Name System (DNS) translates human-readable domain names into the
numerical addresses that are required for almost every transaction on the
Internet~\cite{rfc1035}. Despite this critical role, the original DNS protocol
was not designed with any mechanism for verifying that a response actually
originates from the authoritative source or that it has not been altered in
transit~\cite{rfc4033}. The DNS Security Extensions (DNSSEC) were introduced to
close this gap by adding data origin authentication and data integrity
protection to DNS responses through digital signatures, without, however,
providing confidentiality~\cite{rfc4033,rfc4034,rfc4035}. More than two decades
after its initial specification, DNSSEC is considered the best current practice
for authenticating DNS data~\cite{rfc9364}, even though deployment remains
uneven: recent estimates suggest that fewer than ten percent of website domain
names are signed and that only around a third of queries to recursive resolvers
are actually validated~\cite{rfc9364}.

The cryptographic foundation of DNSSEC, however, is now under long-term threat.
The signature algorithms that are most widely deployed today, in particular RSA
and the Elliptic Curve Digital Signature Algorithm
(ECDSA)~\cite{rfc6605,rfc8624}, derive their security from the computational
hardness of integer factorization and the discrete logarithm problem. A
sufficiently large, cryptographically relevant quantum computer would be able
to solve both of these problems efficiently using Shor's
algorithm~\cite{shor}, which would allow an attacker to forge DNSSEC signatures
and therefore to present spoofed DNS data as authentic.\footnote{This is a
long-term, not an immediate, threat: at the time of writing, no
cryptographically relevant quantum computer is publicly known to exist.
Nevertheless, migration planning for long-lived infrastructure such as DNS is
widely recommended well in advance of such a scenario.} To address this risk,
the U.S. National Institute of Standards and Technology (NIST) has standardized
the Module-Lattice-Based Digital Signature Algorithm (ML-DSA) in FIPS~204, a
signature scheme whose security relies on the hardness of lattice problems and
that is believed to remain secure even against an adversary equipped with a
quantum computer~\cite{fips204}.

Post-quantum cryptography (PQC) schemes such as ML-DSA, however, do not come
without cost. The smallest standardized ML-DSA parameter set, ML-DSA-44,
produces signatures of 2,420~octets and public keys of
1,312~octets~\cite{fips204,mldsadraft}, which is one to two orders of magnitude
larger than a typical ECDSA signature. For comparison, an ECDSA P-256
signature occupies only 64~bytes, so an ML-DSA-44 signature is almost 38~times
larger~\cite{cloudflareblog}. When such a signature is embedded into a DNSSEC
response together with the DNS header, resource records, and---where
applicable---a DNSKEY record, the resulting message can easily exceed the
conservative 1,232-byte UDP payload size that is commonly recommended in
operational practice today, and it will frequently exceed the roughly
1,500-byte Maximum Transmission Unit (MTU) of a typical Ethernet
path~\cite{rfc6891,cloudflaremtu}. Because DNS was originally designed around a
512-byte UDP message size~\cite{rfc1035} and later extended via EDNS(0) to
negotiate larger UDP payloads~\cite{rfc6891}, oversized post-quantum DNSSEC
responses risk falling back to one of two imperfect mechanisms: IP-layer
fragmentation, which is unreliable because many firewalls and other middleboxes
silently drop fragmented packets~\cite{rfc7766,fragpoison}, or truncation of
the UDP response followed by a retry over TCP, which guarantees delivery but
adds at least one additional round trip and extra server-side
state~\cite{rfc7766}. A drop-in replacement of classical signatures with
quantum-safe ones is expected to increase DNSSEC query resolution times by at
least a factor of two: a response carrying large post-quantum signatures is
likely to be truncated, which wastes the initial UDP round trip, and the
subsequent retry over TCP incurs at least two further round trips because of
the three-way handshake~\cite{turbodns}.

The DNS operator community has already begun to experiment with mitigations for
this problem. Cloudflare has deployed experimental ML-DSA-44 validation on its
public 1.1.1.1 resolver, relying on TCP fallback whenever a response becomes
oversized~\cite{cloudflareblog}, and an open-source extension of the CoreDNS
authoritative server has demonstrated on-the-fly signing with several PQC
signature families, including ML-DSA~\cite{corednspqc,coredns}. At the same
time, the research literature has proposed alternative transport mechanisms
that try to avoid the drawbacks of both IP fragmentation and plain TCP
fallback. Application-layer Request-based Resource Record Fragmentation (ARRF)
fragments oversized DNS resource records at the DNS protocol layer itself
rather than relying on the IP layer, with the resolver requesting any missing
fragments explicitly~\cite{arrf}. DNS over QUIC (DoQ) instead moves the entire
DNS exchange onto a dedicated, encrypted QUIC connection, which has its own
packetization and loss-recovery behavior and therefore sidesteps classical IP
fragmentation entirely~\cite{rfc9250}.

Although each of these mechanisms---plain TCP fallback, ARRF, and DoQ---has
been described or prototyped individually, a direct, controlled, and
reproducible comparison of their latency, bandwidth, and reliability behavior
specifically under the conditions created by post-quantum DNSSEC signatures is
still largely missing from the public literature. This gap motivates the
present project. Using ECDSA-based DNSSEC as a classical baseline and an
ML-DSA-44-based DNSSEC deployment as the post-quantum case, this project
evaluates DNS resolution under simulated network conditions with varying
latency and packet loss, measuring query latency, resource utilization,
bandwidth consumption, packet sizes, and the frequency of TCP fallback. Time
permitting, the evaluation is extended to ARRF-based fragmentation and to DNS
over QUIC, in order to assess which combination of signature algorithm and
transport mechanism currently represents the most realistic path toward a
quantum-resistant DNSSEC deployment.

Concretely, the project is guided by the following research questions, which
will be answered in the practical part (Project -- Part~2):
\begin{itemize}
\item[\textbf{RQ1:}] How much do ML-DSA-44 signatures increase DNSSEC response
sizes, packet sizes, and bandwidth consumption, and how often do they trigger
truncation and TCP fallback, compared with the ECDSA baseline?
\item[\textbf{RQ2:}] How do the end-to-end query latency and resource
utilization of post-quantum DNSSEC with TCP fallback change as path latency and
packet loss increase?
\item[\textbf{RQ3:}] (Time permitting.) To what extent can ARRF or DNS over
QUIC reduce the overhead identified in RQ1 and RQ2?
\end{itemize}

The remainder of this (theoretical) part of the report is structured as
follows. Section~\ref{sec:foundations} introduces the theoretical foundations
that are necessary to understand the problem: the DNS and DNSSEC protocols, the
classical and post-quantum signature algorithms relevant to this project, and
the mechanisms that govern DNS message size and transport selection.
Section~\ref{sec:related} positions this project with respect to the existing
state of the art and related work. Section~\ref{sec:conclusion} closes the
theoretical part with a brief summary and a transition toward the practical
implementation and evaluation that are reported in Part~2 of this project.

\begin{thebibliography}{00}

%BIB-BEGIN
\bibitem{rfc1035} P. Mockapetris, ``Domain names---Implementation and specification,'' RFC~1035, STD~13, Internet Engineering Task Force (IETF), Nov. 1987. [Online]. Available: \url{https://www.rfc-editor.org/rfc/rfc1035}. [Accessed: Oct. 20, 2026].

\bibitem{rfc4033} R. Arends, R. Austein, M. Larson, D. Massey, and S. Rose, ``DNS security introduction and requirements,'' RFC~4033, Internet Engineering Task Force (IETF), Mar. 2005. [Online]. Available: \url{https://www.rfc-editor.org/rfc/rfc4033}. [Accessed: Oct. 20, 2026].

\bibitem{rfc4034} R. Arends, R. Austein, M. Larson, D. Massey, and S. Rose, ``Resource records for the DNS security extensions,'' RFC~4034, Internet Engineering Task Force (IETF), Mar. 2005. [Online]. Available: \url{https://www.rfc-editor.org/rfc/rfc4034}. [Accessed: Oct. 20, 2026].

\bibitem{rfc4035} R. Arends, R. Austein, M. Larson, D. Massey, and S. Rose, ``Protocol modifications for the DNS security extensions,'' RFC~4035, Internet Engineering Task Force (IETF), Mar. 2005. [Online]. Available: \url{https://www.rfc-editor.org/rfc/rfc4035}. [Accessed: Oct. 20, 2026].

\bibitem{rfc9364} P. Hoffman, ``DNS security extensions (DNSSEC),'' RFC~9364, BCP~237, Internet Engineering Task Force (IETF), Feb. 2023. [Online]. Available: \url{https://www.rfc-editor.org/rfc/rfc9364}. [Accessed: Oct. 20, 2026].

\bibitem{rfc6605} P. Hoffman and W. C. A. Wijngaards, ``Elliptic curve digital signature algorithm (DSA) for DNSSEC,'' RFC~6605, Internet Engineering Task Force (IETF), Apr. 2012. [Online]. Available: \url{https://www.rfc-editor.org/rfc/rfc6605}. [Accessed: Oct. 20, 2026].

\bibitem{rfc8624} P. Wouters and O. Sury, ``Algorithm implementation requirements and usage guidance for DNSSEC,'' RFC~8624, Internet Engineering Task Force (IETF), Jun. 2019. [Online]. Available: \url{https://www.rfc-editor.org/rfc/rfc8624}. [Accessed: Oct. 20, 2026].

\bibitem{shor} P. W. Shor, ``Polynomial-time algorithms for prime factorization and discrete logarithms on a quantum computer,'' \emph{SIAM J. Comput.}, vol.~26, no.~5, pp.~1484--1509, 1997, doi: 10.1137/S0097539795293172.

\bibitem{fips204} National Institute of Standards and Technology, \emph{Module-Lattice-Based Digital Signature Standard}, FIPS~204, Aug. 2024, doi: 10.6028/NIST.FIPS.204. [Online]. Available: \url{https://csrc.nist.gov/pubs/fips/204/final}. [Accessed: Oct. 20, 2026].

\bibitem{mldsadraft} B. E. Westerbaan and S. Schmieg, ``Module-lattice digital signature algorithm for DNSSEC,'' IETF Internet-Draft draft-westerbaan-dnssec-mldsa-04 (work in progress), Aug. 2026. [Online]. Available: \url{https://datatracker.ietf.org/doc/draft-westerbaan-dnssec-mldsa/}. [Accessed: Oct. 20, 2026].

\bibitem{cloudflareblog} S. Neuteboom and B. Westerbaan, ``1.1.1.1 now supports post-quantum DNSSEC, all 2,420 bytes of it,'' Cloudflare Blog, Sep. 10, 2026. [Online]. Available: \url{https://blog.cloudflare.com/post-quantum-dnssec-1111/}. [Accessed: Oct. 2, 2026].

\bibitem{rfc6891} J. Damas, M. Graff, and P. Vixie, ``Extension mechanisms for DNS (EDNS(0)),'' RFC~6891, STD~75, Internet Engineering Task Force (IETF), Apr. 2013. [Online]. Available: \url{https://www.rfc-editor.org/rfc/rfc6891}. [Accessed: Oct. 20, 2026].

\bibitem{cloudflaremtu} Cloudflare, ``Maximum transmission unit and maximum segment size,'' Cloudflare WAN Documentation, Aug. 24, 2026. [Online]. Available: \url{https://developers.cloudflare.com/cloudflare-wan/reference/mtu-mss/}. [Accessed: Oct. 2, 2026].

\bibitem{rfc7766} J. Dickinson, S. Dickinson, R. Bellis, A. Mankin, and D. Wessels, ``DNS transport over TCP---Implementation requirements,'' RFC~7766, Internet Engineering Task Force (IETF), Mar. 2016. [Online]. Available: \url{https://www.rfc-editor.org/rfc/rfc7766}. [Accessed: Oct. 20, 2026].

\bibitem{fragpoison} A. Herzberg and H. Shulman, ``Fragmentation considered poisonous, or: One-domain-to-rule-them-all.org,'' in \emph{Proc. 2013 IEEE Conf. Communications and Network Security (CNS)}, 2013, pp.~224--232, doi: 10.1109/CNS.2013.6682711.

\bibitem{turbodns} A. S. Rawat and M. P. Jhanwar, ``Post-quantum DNSSEC with faster TCP fallbacks,'' in \emph{Progress in Cryptology -- INDOCRYPT 2024}, Lecture Notes in Computer Science. Cham, Switzerland: Springer, 2024, pp.~212--236, doi: 10.1007/978-3-031-80311-6\_11.

\bibitem{corednspqc} J. Gento Suela, J. Blanco-Romero, F. Almenares Mendoza, and D. D\'iaz-S\'anchez, ``Implementing and evaluating post-quantum DNSSEC in CoreDNS,'' in \emph{Proc. 2025 Int. Conf. Modeling, Analysis and Simulation of Wireless and Mobile Systems (MSWiM)}, 2025, pp.~420--427, doi: 10.1109/MSWIM67937.2025.11308970.

\bibitem{coredns} CoreDNS, ``CoreDNS,'' GitHub repository. [Online]. Available: \url{https://github.com/coredns/coredns}. [Accessed: Oct. 2, 2026].

\bibitem{arrf} J. Goertzen and D. Stebila, ``Post-quantum signatures in DNSSEC via request-based fragmentation,'' in \emph{Proc. 14th Int. Conf. Post-Quantum Cryptography (PQCrypto 2023)}, Lecture Notes in Computer Science, vol.~14154, 2023, pp.~535--564, doi: 10.1007/978-3-031-40003-2\_20.

\bibitem{rfc9250} C. Huitema, S. Dickinson, and A. Mankin, ``DNS over dedicated QUIC connections,'' RFC~9250, Internet Engineering Task Force (IETF), May 2022. [Online]. Available: \url{https://www.rfc-editor.org/rfc/rfc9250}. [Accessed: Oct. 2, 2026].
%BIB-END

\end{thebibliography}

\end{document}
